\section*{अध्याय \ref{ch:b1:extent-q-and-k} --- रासायनिक निकाय का वर्णन: प्रगति, सक्रियताएँ, Q और K}

\begin{solution}{exo:b1:extent-q-and-k:1}
\omterm{def:b1:extent-q-and-k:variables}{गहन}: ताप, \omterm{def:b1:extent-q-and-k:composition}{मोलर सांद्रता}, घनत्व, दाब, \omterm{def:b1:extent-q-and-k:composition}{मोल
अंश}। \omterm{def:b1:extent-q-and-k:variables}{विस्तीर्ण}: आयतन, पदार्थ की मात्रा, द्रव्यमान।
\end{solution}

\begin{solution}{exo:b1:extent-q-and-k:2}
$p(\ce{N2}) = 0.7808 \times 1.013 = \qty{0.791}{bar}$, $p(\ce{O2}) =
0.2095 \times 1.013 = \qty{0.212}{bar}$, $p(\ce{Ar}) = 0.0093 \times 1.013
= \qty{0.0094}{bar}$।
\end{solution}

\begin{solution}{exo:b1:extent-q-and-k:3}
$n(\ce{H2}) = 3.0 - 2\xi$, $n(\ce{O2}) = 2.0 - \xi$, $n(\ce{H2O}) =
2\xi$। $\xi_{\max} = \min(3.0/2,\ 2.0/1) = \qty{1.5}{mol}$: डाइहाइड्रोजन
सीमांत है। अंतिम अवस्था: \ce{H2} 0, \ce{O2} \qty{0.5}{mol}, \ce{H2O}
\qty{3.0}{mol}।
\end{solution}

\begin{solution}{exo:b1:extent-q-and-k:4}
(a) $Q = p_{\ce{CO2}}/p^\circ$। (b) $Q = (c^\circ)^2/([\ce{Ag+}][\ce{Cl-}])$।
(c) $Q = (p_{\ce{NH3}}/p^\circ)^2/[(p_{\ce{N2}}/p^\circ)(p_{\ce{H2}}/p^\circ)^3]$।
(d) $Q = [\ce{NH4+}][\ce{OH-}]/([\ce{NH3}]\,c^\circ)$, क्योंकि पानी की \omterm{def:b1:extent-q-and-k:activity}{सक्रियता} 1 है।
\end{solution}

\begin{solution}{exo:b1:extent-q-and-k:5}
$\mathrm{A \rightleftharpoons C}$ योग है: $K = 10 \times 0.5 = 5$। $\mathrm{C \rightleftharpoons A}$: $1/5 =
0.2$। $\mathrm{2\,A \rightleftharpoons 2\,C}$: $5^2 = 25$।
\end{solution}

\begin{solution}{exo:b1:extent-q-and-k:6}
$Q = 5.0 \times 10^{-5} < K^\circ$: अग्र। $Q = 0.10 > K^\circ$: प्रतीप।
केवल अभिकारक: $Q = 0 < K^\circ$, अग्र।
\end{solution}

\begin{solution}{exo:b1:extent-q-and-k:7}
$n(\ce{A2}) = 1 - \xi$, $n(\ce{A}) = 2\xi$, कुल $1 + \xi$। $p =
p^\circ$ के साथ, $Q = \dfrac{(2\xi/(1+\xi))^2}{(1-\xi)/(1+\xi)} = \dfrac{4\xi^2}{1 -
\xi^2}$। $Q = 0.25$ से $4.25\,\xi^2 = 0.25$, $\xi = \qty{0.243}{mol}$।
\end{solution}

\begin{solution}{exo:b1:extent-q-and-k:8}
$x = [\ce{C}]$ लेकर: $x/(0.10 - x)^2 = 100$, इसलिए $100x^2 - 21x + 1 = 0$ और
$x = (21 - \sqrt{41})/200 = \qty{0.0730}{mol/L}$ (दूसरा मूल
0.10 से अधिक है)। $[\ce{A}] = [\ce{B}] = \qty{0.0270}{mol/L}$, $[\ce{C}] =
\qty{0.0730}{mol/L}$, $\tau = 0.73$।
\end{solution}

\begin{solution}{exo:b1:extent-q-and-k:9}
हर एक के $n$ से: $Q = \xi^2/(n - \xi)^2 = K^\circ$, इसलिए $\xi/(n - \xi) =
\sqrt{K^\circ}$ और $\tau = \xi/n = \sqrt{K^\circ}/(1 + \sqrt{K^\circ})$।
$\tau = 0.99$ के लिए $\sqrt{K^\circ} = 99$ चाहिए, $K^\circ \approx \num{9.8e3}$:
लगभग $10^4$, \omterm{def:b1:extent-q-and-k:quantitative}{मात्रात्मक अभिक्रिया} की सामान्य सीमा।
\end{solution}

\begin{solution}{exo:b1:extent-q-and-k:10}
साम्य पर $p_{\ce{CO2}} = 0.20\,p^\circ$, अर्थात् $n(\ce{CO2}) =
pV/RT = 0.20 \times 10^5 \times 0.0100/(8.314 \times 1100) =
\qty{0.0219}{mol}$। कार्बोनेट के \qty{0.010}{mol} से यह प्राप्त नहीं
हो सकता: पूरा कार्बोनेट विघटित हो जाता है, $p_{\ce{CO2}} = 0.010 \times 8.314 \times
1100/0.0100 = \qty{9.15e3}{Pa} = \qty{0.091}{bar}$, $Q < K^\circ$, और
अंतिम अवस्था (\qty{0.010}{mol} \ce{CaO}, \qty{0.010}{mol} \ce{CO2}, कोई
\ce{CaCO3} नहीं) साम्य नहीं है। \qty{0.050}{mol} के साथ: साम्य,
\qty{0.0219}{mol} \ce{CO2} और \ce{CaO}, और \qty{0.0281}{mol} \ce{CaCO3} शेष।
\end{solution}

\begin{solution}{exo:b1:extent-q-and-k:11}
साम्य पर $n(\ce{B}) = 2\,n(\ce{A})$ और $n(\ce{C}) = 3\,n(\ce{A})$
(आयतन कट जाता है)। तब $6\,n(\ce{A}) = 1.0$: $n(\ce{A}) =
\qty{0.167}{mol}$, $n(\ce{B}) = \qty{0.333}{mol}$, $n(\ce{C}) =
\qty{0.500}{mol}$।
\end{solution}

\begin{solution}{exo:b1:extent-q-and-k:12}
बराबर मात्राएँ: $\xi^2/(0.50 - \xi)^2 = 4$, $\xi/(0.50 - \xi) = 2$,
$\xi = \qty{0.333}{mol}$, \omterm{def:b1:extent-q-and-k:fractional-extent}{लब्धि} $0.333/0.50 = 67\,\%$। ऐल्कोहॉल के \qty{1.0}{mol}
के साथ: $\xi^2 = 4(0.50 - \xi)(1.0 - \xi)$, $3\xi^2 - 6\xi + 2 = 0$,
$\xi = (6 - \sqrt{12})/6 = \qty{0.423}{mol}$, \omterm{def:b1:extent-q-and-k:fractional-extent}{लब्धि} 85\,\%: एक अभिकारक का
आधिक्य दूसरे के उत्पाद की \omterm{def:b1:extent-q-and-k:fractional-extent}{लब्धि} बढ़ा देता है।
\end{solution}

\begin{solution}{pb:b1:extent-q-and-k:1}
\textbf{1.} $x(\ce{CH4}) = 1.00/4.00 = 0.250$; $x(\ce{H2O}) = 0.750$।
\textbf{2.} \qty{7.5}{bar} और \qty{22.5}{bar}।
\textbf{3.} द्रव्यमान 16.04 और \qty{54.06}{g}: $w(\ce{CH4}) = 0.229$,
$w(\ce{H2O}) = 0.771$।
\textbf{4.} \omterm{def:b1:extent-q-and-k:composition}{मोल अंश}, \omterm{def:b1:extent-q-and-k:composition}{द्रव्यमान अंश}, \omterm{def:b1:extent-q-and-k:composition}{आंशिक दाब} (और
कुल दाब); मात्राएँ \omterm{def:b1:extent-q-and-k:variables}{विस्तीर्ण} हैं।
\textbf{5.} $a = p_i/p^\circ$: मेथेन के लिए 7.5, भाप के लिए 22.5।
\textbf{6.} \ce{CH4} $1.00 - \xi$, \ce{H2O} $3.00 - \xi$, \ce{CO} $\xi$,
\ce{H2} $3\xi$, कुल $4.00 + 2\xi$।
\textbf{7.} $\xi_{\max} = \qty{1.00}{mol}$; मेथेन सीमांत है।
\textbf{8.} \ce{CH4} 0, \ce{H2O} \qty{2.00}{mol}, \ce{CO} \qty{1.00}{mol},
\ce{H2} \qty{3.00}{mol}।
\textbf{9.} \qty{6.00}{mol}, फ़ीड के \qty{4.00}{mol} के मुक़ाबले: अभिक्रिया
दो अणुओं से चार अणु बनाती है।
\textbf{10.} \ce{H2O} 0.333, \ce{CO} 0.167, \ce{H2} 0.500।
\textbf{11.} ताकि पूरी मेथेन अभिक्रिया कर ले और सृति अभिक्रिया के लिए
भाप बचे, जिसका भागफल वह घटाती है।
\textbf{12.} \ce{CO} $1 - \xi$, \ce{H2O} $2 - \xi$, \ce{CO2} $\xi$, \ce{H2}
$3 + \xi$, हर प्रगति पर कुल 6.00।
\textbf{13.} $Q = \dfrac{(p_{\ce{CO2}}/p^\circ)(p_{\ce{H2}}/p^\circ)}
{(p_{\ce{CO}}/p^\circ)(p_{\ce{H2O}}/p^\circ)}$, जहाँ $p_i = (n_i/6)p$: गुणक
$p/(6p^\circ)$ अंश और हर के बीच कट जाते हैं (हर ओर दो गैस
अणु), और दिया गया व्यंजक बचता है।
\textbf{14.} $\xi = 0$: $Q = 0 < 4.0$, अग्र।
\textbf{15.} $\xi(3 + \xi) = 4(1 - \xi)(2 - \xi)$, अर्थात् $3\xi^2 - 15\xi +
8 = 0$: $\xi = (15 - \sqrt{129})/6 = \qty{0.607}{mol}$; दूसरा मूल,
4.39, $\xi_{\max} = 1$ से अधिक है।
\textbf{16.} \ce{CO} \qty{0.393}{mol}, \ce{H2O} \qty{1.393}{mol},
\ce{CO2} \qty{0.607}{mol}, \ce{H2} \qty{3.607}{mol}।
\textbf{17.} $\tau = 0.607/1.00 = 0.607$।
\textbf{18.} $\xi(3 + \xi) = 4(1 - \xi)(5 - \xi)$, $3\xi^2 - 27\xi + 20 =
0$, $\xi = \qty{0.814}{mol}$: अधिक भाप, अधिक रूपांतरण।
\textbf{19.} $K^\circ = (0.99 \times 3.99)/(0.01 \times 1.01) = 391$।
\textbf{20.} $Q(\xi)$ वर्धमान है, इसलिए बड़ा $K^\circ$ बड़ी प्रगति पर
मिलता है। इस अभिक्रिया का स्थिरांक कम ताप पर बड़ा होता है, जहाँ अभिक्रिया
धीमी होती है: पहला रिएक्टर गर्म और तेज़ चलता है, दूसरा, ठंडा, रूपांतरण पूरा करता है
($K^\circ$ की ताप-निर्भरता द्वितीय वर्ष के खंड में दी गई है)।
\textbf{21.} $3.607/(6.00 - 1.393) = 3.607/4.607 = 0.783$।
\textbf{22.} \ce{CO} का \qty{0.393}{mol} और \ce{CO2} का \qty{0.607}{mol}।
\textbf{23.} $x(\ce{H2}) = 3.607/6.00 = \textbf{0.60}$: सृति रिएक्टर से निकलने वाले हर
दस अणुओं में छह हाइड्रोजन के हैं।
\end{solution}
